% ══════════════════════════════════════════════════════════════════════════════
%  EXPERIMENT 3 — QUADRAT SURVEY                         (Lesson 5, T3 Week 1)
% ══════════════════════════════════════════════════════════════════════════════
\experimentdiv{Quadrat Survey}{Sampling with quadrats \textbullet{} Averages from repeat samples \textbullet{} Comparing two sites}{5}\label{sec:3}
\lessoncodes{VC2S8U04 \textbullet{} VC2S8I01, I02, I04, I06}

\vocabhead
\begin{ctab}
\begin{tabular}{|SW{4.2cm}|W{11.6cm}|}
  \hline\hd & \thead{Word} & \thead{Definition (copy from the slide)} \\ \hline
  & \scaf{\textbf{Quadrat}} & \scaf{A square frame, often 1\,m by 1\,m, used to sample the organisms in one small area.} \\ \hline
  & \textbf{Biodiversity} & \\ \hline
  & \textbf{Species richness} & \\ \hline
  & \textbf{Sample} & \\ \hline
  & \textbf{Random sampling} & \\ \hline
\end{tabular}
\end{ctab}

\begin{lipanel}
\inquiryli{Collect data with quadrats, average repeat samples and compare two sites in a column graph.}
\understandli{Compare the biodiversity of two habitats and explain the difference using abiotic and biotic factors.}
\sctitle
\begin{itemize}[leftmargin=1.4em, itemsep=0pt, topsep=1pt]
  \item I can place quadrats randomly and count what is inside.
  \item I can calculate an average from three quadrats.
  \item I can explain why one site has more species than the other.
\end{itemize}
\end{lipanel}

\begin{skillsfocus}
\focus{QP1}{Write a hypothesis about the two sites.}
\focus{PC1}{Follow the sampling method exactly.}
\focus{PC6}{Count every species in each quadrat.}
\focus{PA1}{Record counts in a THUD table.}
\focus{PA2}{Draw a column graph of the averages.}
\focus{PA3}{Calculate averages.}
\focus{EV3}{Explain how sampling can give errors.}
\end{skillsfocus}

% ─── TECHNIQUE ────────────────────────────────────────────────────────────────
\exppart{Technique spotlight}{Blue}

\begin{technique}{Why we sample}{PC1}
You cannot count every plant and animal on the whole oval, so you count a few small areas (\textbf{samples}) and use them to stand for the whole site. Two rules keep a sample fair:
\begin{itemize}[leftmargin=1.4em, itemsep=1pt, topsep=3pt]
  \item \textbf{Random:} you do not choose the interesting spots. Stand still, close your eyes and toss the quadrat gently over your shoulder.
  \item \textbf{Repeated:} one quadrat could land on an unusual patch. Take at least three and average them.
\end{itemize}
\end{technique}

\begin{technique}{Counting species richness}{PC6}
\begin{minipage}[c]{0.6\linewidth}
\textbf{Species richness} is the number of \emph{different} kinds of living things, not the total number of living things. Ten clover plants and one ant is a richness of 2. You do not need names: ``broad grass'', ``thin grass'', ``small black ant'' are fine, as long as you use the same name at both sites.
\end{minipage}\hfill
\begin{minipage}[c]{0.36\linewidth}\centering
\begin{tikzpicture}[scale=0.55]
  \draw[Ink, line width=1.6pt] (0,0) rectangle (4,4);
  \foreach \x/\y in {0.6/0.7,1.4/2.9,2.6/1.2,3.3/3.2,0.8/3.4,2.0/2.2}{\draw[Ink, line width=0.8pt] (\x,\y) -- (\x,\y+0.45) (\x,\y) -- (\x-0.2,\y+0.35) (\x,\y) -- (\x+0.2,\y+0.35);}
  \foreach \x/\y in {1.2/1.0,3.2/0.6,2.8/2.7}{\draw[Ink, fill=white, line width=0.8pt] (\x,\y) circle (0.22); \foreach \a in {0,72,...,288}{\draw[Ink] (\x,\y) ++(\a:0.22) -- ++(\a:0.14);}}
  \fill[Ink] (3.6,1.8) ellipse (0.12 and 0.07);
  \node[font=\scriptsize, anchor=north] at (2,-0.1) {1\,m};
\end{tikzpicture}
\end{minipage}
\end{technique}

\begin{workedex}{Richness and the average}
Three quadrats on a garden bed had species richness 7, 9 and 8.

\smallskip
Average $= (7 + 9 + 8) \div 3 = 24 \div 3 = 8$ species.
\end{workedex}

\begin{wedo}{Counting richness}
The quadrat in the picture above holds grass, a flowering weed and one ant. What is its species richness?
\answerbox{2cm}
\end{wedo}

% ─── PLAN AND SAFETY / METHOD ────────────────────────────────────────────────
\exppart{Plan, safety and method}{Teal}

\textbf{Aim:} To investigate the effect of changing the site (the mown oval or a garden bed) on species richness.

\task Write a hypothesis. \skr{QP1}
\answerbox{2.5cm}

\needspace{12\baselineskip}
\task Complete the variables table. \skr{PC2}
\begin{ctab}
\begin{tabular}{|SW{4.4cm}|W{10.2cm}|}
  \hline\hd & \thead{Variable} & \thead{In this experiment} \\ \hline
  & \scaf{Independent} & \scaf{The site: mown oval or garden bed} \\ \hline
  & Dependent & \\ \hline
  & Controlled & \\ \hline
  & Controlled & \\ \hline
\end{tabular}
\end{ctab}

\needspace{10\baselineskip}
\task Complete the risk assessment. \skr{PC3}
\begin{ctab}
\begin{tabular}{|SW{4.0cm}|W{5.2cm}|W{5.4cm}|}
  \hline\hd & \thead{Risk (what could go wrong)} & \thead{Prevention} & \thead{Responding (if it happens)} \\ \hline
  & \scaf{Stung by a bee or bull ant} & \scaf{Look before you kneel; do not put hands under logs} & \scaf{Tell the teacher; anyone allergic uses their EpiPen plan} \\ \hline
  & A thrown quadrat hits someone & & \\ \hline
  & Sunburn & & \\ \hline
\end{tabular}
\end{ctab}

\subsubsection*{Equipment}
1\,m $\times$ 1\,m quadrat (or four metre rulers) \textbullet{} hand lens \textbullet{} clipboard and pencil \textbullet{} hat and sunscreen

\subsubsection*{Method}
\begin{enumerate}[itemsep=2pt]
  \item At the first site, stand still with your back to the area, check no one is behind you, and toss the quadrat gently over your shoulder.
  \item Search the quadrat for 3 minutes, including under leaves. Record each different kind of plant and animal, and count how many kinds there are.
  \item Repeat steps 1 and 2 twice more at the same site.
  \item Repeat steps 1 to 3 at the second site.
  \item Calculate the average species richness for each site.
\end{enumerate}

% ─── RESULTS ─────────────────────────────────────────────────────────────────
\exppart{Results}{Amber}

\subsubsection*{Table 3.1: Species richness at two sites \skr{PA1}}
\begin{ctab}
\begin{tabular}{|SW{3.4cm}|C{2.6cm}|C{2.6cm}|C{2.6cm}|C{3.0cm}|}
  \hline\hd & \thead{Site} & \thead{Quadrat 1 (species)} & \thead{Quadrat 2 (species)} & \thead{Quadrat 3 (species)} & \thead{Average (species)} \\ \hline
  & \scaf{Example: path edge} & \scaf{4} & \scaf{6} & \scaf{5} & \scaf{5} \\ \hline
  & Mown oval & & & & \\ \hline
  & Garden bed & & & & \\ \hline
\end{tabular}
\end{ctab}

\task Show your working for one average. \skr{PA3}
\workbox{2.5cm}

\needspace{12\baselineskip}
\task List the kinds of living things you found at each site.
\begin{ctab}
\begin{tabular}{|W{7.4cm}|W{7.4cm}|}
  \hline\thead{Mown oval} & \thead{Garden bed} \\ \hline
  \rule{0pt}{4cm} & \\ \hline
\end{tabular}
\end{ctab}

\needspace{18\baselineskip}
\task Draw a column graph of the average species richness at each site. Give it a title and label both axes with units. \skr{PA2}
\graphgrid{8cm}

% ─── ANALYSIS ────────────────────────────────────────────────────────────────
\exppart{Analysis and evaluation}{Coral}

\task Which site had greater biodiversity? Use numbers from Table 3.1. \skr{PA5}
\answerbox{3.5cm}

\needspace{12\baselineskip}
\task Explain the difference between the sites. Use at least one abiotic factor (such as shade, water or mowing) and one biotic factor.
\answerbox{5cm}

\needspace{12\baselineskip}
\task Give two reasons your averages might not match the true richness of each site, and how a better survey could fix each. \skr{EV3}
\begin{ctab}
\begin{tabular}{|SW{7.2cm}|W{7.8cm}|}
  \hline\hd & \thead{Source of error} & \thead{How to reduce it} \\ \hline
  & \scaf{Only three quadrats per site} & \scaf{Take ten quadrats per site and average them} \\ \hline
  & & \\ \hline
  & & \\ \hline
\end{tabular}
\end{ctab}

\needspace{10\baselineskip}
\task Write a conclusion: the trend, the numbers that back it, and whether your hypothesis was supported. \skr{EV1}
\answerbox{4.5cm}

\begin{yr8box}
\begin{keepext}
\exttask Species richness counts kinds, not how many of each. Garden A has 9 kinds, each with about 10 individuals. Garden B has 9 kinds, but 90\% of the individuals are one grass. Which garden is more diverse? Explain why richness alone does not show this.
\answerbox{5cm}
\end{keepext}
\end{yr8box}

% ─── SKILL CHECK ─────────────────────────────────────────────────────────────
\exppart{Skill check}{NavyBlue}

\begin{skillcheck}
\scrow{QP1}{I can write a hypothesis that names both variables.}
\scrow{PC1}{I followed the sampling method exactly, including random placement.}
\scrow{PC6}{I counted every kind of living thing in each quadrat.}
\scrow{PA1}{My table has a title, headings, units and matching decimals.}
\scrow{PA2}{I drew a column graph with a title and labelled axes.}
\scrow{PA3}{I calculated averages correctly.}
\scrow{EV3}{I can explain how sampling can give errors.}
\end{skillcheck}

\task Which skill will you focus on in Experiment 4? Explain what you will do differently.
\answerbox{5cm}
